\subsection{SW12: the numerical interfaces for the actual kept blocks}
\label{sec:segment-numerics}

\status{The literal standard hypotheses imply $\mathrm{BLk}(0)$,
$o_+\le a_+$, and the source total-order bound with $C_t\ge30$.
The uniform supply of those hypotheses and the comparison $t_0\le t_I$
remain separate.}

Assume $\mathrm{Std}(C_0,C_\delta,C_{E1},C_{E2},1,X)$.
Write $k=m+1$, $\lambda=X.\mathrm{La}(C_0)$, $\delta=X.\mathrm{de}$,
$\beta=\mathrm{be}(X,C_0)$ and $\epsilon=N_1\lambda/k$.
The clauses used below give
\[
 \begin{gathered}
 p\ge5,\quad k\ge5,\quad p\ge k^2,\quad
 40\epsilon\le\delta<1/6,\quad
 0\le\beta\le\delta/2-18\epsilon,\\
 \lambda/k\le1/4,\quad \epsilon\ge1/k,\quad
 b'\le a/p,\quad k^3Q\le b'.
 \end{gathered}
\]
In one actual segment, $G\ge2$ is even and its interior distances
$0<d_0<\cdots<d_{s-1}<G$ are odd. Put $P_i=p^{G-d_i}$,
$g_i=b_T(\mathrm{pos}(j,i))$, $A_0=\lfloor(a-E_0)/p^G\rfloor$,
and let $A_i$ be the coarse degree after the $k$ compatible digits.
Thus $g_i-kP_i<A_iP_i\le g_i$ and
$b'(1-\beta)\le g_i\le b'(1+\beta)\le2b'$.
The split is $i<\sigma\iff p^{d_i}<a/b'$.
For the effective order $e_i=w_i+\mathrm{sg}_i$ before the split and
$e_i=o_i$ after it, the signed integer cap is
$\mathrm{cap}_l=A_0+\sum_{i<l}(A_i-e_i)$.
No positivity of $A_0$ or of a cap is assumed. All cap comparisons below
are upper bounds; the nonnegative effective orders can therefore be
discarded when obtaining them. The remaining window inequalities are
literal clauses of $\mathrm{Std}$.

\paragraph{The late cutoff.}
Odd spacing and $E_0\ge0$ give
\[
 P_i\mathrm{cap}_i
 \le\frac a{p^{d_i}}+2b'\sum_{h<i}p^{d_h-d_i}
 \le\frac a{p^{d_i}}+\frac{2b'}p.
 \tag{SW.12a}
\]
Indeed the geometric sum is at most $(p^2-1)^{-1}\le p^{-1}$.
Also $2/p\le\epsilon$ and $kP_i\le kQ\le\epsilon b'$.
For $i\ge\sigma$, the excluded-band clause of $\mathrm{Std}$ strengthens
the failed split test to $p^{d_i}>(a/b')(1+2\delta)$.
Set $t=a/p^{d_i}$, $u=1-\delta/2+17\epsilon$ and
$v=1-\delta/4-\epsilon$. Then
\[
 t<\frac{b'}{1+2\delta},\qquad
 P_i\mathrm{cap}_i\le t+\epsilon b',\qquad A_iP_i>b'u.
\]
Both $u,v$ are positive. The explicit remainders
\[
\begin{aligned}
 uv-\epsilon-(1-\delta/2)(1-\delta/4)
   &=\epsilon(15-15\delta/4-17\epsilon)\ge0,\\
 (1+2\delta)(1-\delta/2)(1-\delta/4)-1
   &=\delta(5/4-11\delta/8+\delta^2/4)\ge0
\end{aligned}
\]
show $P_i\mathrm{cap}_i<A_iP_i v$.
Adding $(\lambda/k)A_iP_i\le\epsilon A_iP_i$ proves
\[
 \mathrm{cap}_i+(\lambda/k)A_i\le A_i(1-\delta/4).
 \tag{SW.12b}
\]
The kept-block upper bound for $o_i$ now gives the strengthened late
clause of $\mathrm{BLk}(0)$; all its early clauses are already in
the definition of a kept block.

\paragraph{The terminal extra order.}
The distinct positive exponents $G-d_i$ imply
$\sum_iP_i^{-1}\le(p-1)^{-1}\le5/(4p)$.
Consequently $\mathrm{cap}_s\le a/p^2+3b'/p$ and
$\sum_i2a/P_i\le4a/p$. The literal kept terminal bound is
\[
 o_+\le\mathrm{cap}_s+
  (\lambda/k+4/p^2)\sum_i2a/P_i+\lambda b_+/k,
 \qquad b_+\le2b'.
\]
It follows that
\[
\begin{aligned}
 o_+&\le L_{\rm cap}:=
  \frac a{p^2}+\frac{3b'}p+
  \left(\frac\lambda k+\frac4{p^2}\right)\frac{4a}p
  +\frac{2\lambda b'}k\\
 &\le\left[\frac1{25}+\frac3{25}
       +\left(\frac14+\frac4{25}\right)\frac45+\frac1{10}\right]a
   =\frac{147}{250}a\le a_+.
\end{aligned}
\tag{SW.12c}
\]
Here $L_{\rm cap}$ is exactly \texttt{localOrderCap}, and the terminal
window gives $a_+\ge a$. This proves the remaining $\mathrm{BLk}(0)$
clauses and the conclusion required by \texttt{OPlusAt\_of} whenever
its explicit hypothesis supplying $\mathrm{Std}$ is available.

\paragraph{The full order, with its minimum term.}
Put
\[
 R=\frac a{p^2}+\frac{b'}p+\frac{\lambda a}{pk}
     +\frac{\lambda b'}k+kQ,
 \qquad M_i=\min\{g_i,a_{\rm start}/p^{d_i}\}.
\]
The early excluded-band inequality gives
$g_i\le a_{\rm start}/p^{d_i}$; the late one gives
$a/p^{d_i}\le M_i$.
Multiplying the respective kept inequalities
$e_i\le A_i(1+4/p^2)+(\lambda/k)\mathrm{cap}_i$ and
$e_i\le\mathrm{cap}_i+(\lambda/k)A_i$ by $P_i$, and using
(SW.12a), therefore proves
\[
 P_i e_i\le M_i+T,\qquad
 T=\frac{2b'}p+\frac{2\lambda b'}k+\frac{8b'}{p^2}
       +\frac{\lambda a}{pk}+\frac{2\lambda b'}{pk}.
\]
The digit contribution is at most $2kQ$ per frame. Using $P_i\ge p$
in the terminal bound, and $1/p\le1/5$, gives the following coefficients
of the five summands of $R$ in the complete error:
\[
 (1+2s),\quad6s,\quad3s,\quad4s+2,\quad2s.
\]
Each is at most $30(s+1)$. Thus the actual natural-number order obeys
\[
 \mathrm{blockOrd}_j\le\sum_iM_{j,i}+30(s_j+1)R.
 \tag{SW.12d}
\]
The cyclic segment interiors are disjoint; each contributes its selected
odd points and one distinct even start. Hence
$\sum_j(s_j+1)\le|TT|$. Summing (SW.12d), for $C_t\ge30$ and $c\ge0$,
\[
 \sum_j\mathrm{blockOrd}_j
 \le\sum_{j,i}M_{j,i}+C_t|TT|(R+Qc)=X.t_0(C_0,C_t).
 \tag{SW.12e}
\]
This is the literal \texttt{ord\_all} endpoint. Its proof also works for
$\mathrm{Std}$ with arbitrary $H$, since this last order estimate uses
neither of the clauses containing $H$.
